\begin{tomtat}
\begin{itemize}
  \item Bài báo của chương này chứng minh định lý và không đo gì. Đối tượng nó
  chặn là sai số giữa đầu ra của lời giải thích và đầu ra của mô hình, tức độ
  khớp; rà toàn văn 65 trang thì chữ \enquote{faithfulness} và chữ
  \enquote{ground truth} không xuất hiện lần nào, \enquote{attribution} xuất
  hiện một lần ở phần mở đầu, còn LIME và SHAP xuất hiện một lần, trong cùng
  một câu ở trang 39.
  \item Khung của bài đo độ đơn giản của một lời giải thích bằng độ phức tạp
  Kolmogorov, tức độ dài chương trình ngắn nhất sinh ra nó, và biến
  \enquote{dễ diễn giải} thành một hạn mức $b$ bit. Từ đó ra hai hàm đối ngẫu
  nhau: $\varepsilon_f(b)$, sai số nhỏ nhất trong ngân sách $b$, và
  $\kappa_f(\delta)$, ngân sách nhỏ nhất đạt sai số $\delta$.
  Hình~\ref{fig:ch13-mien-kha-thi} vẽ cả hai cùng một lúc.
  \item Cái trần là định lý~\ref{thm:ch13-khoang-cach-do-phuc-tap}, chứng minh
  dài bốn dòng và không giả thiết gì về mô hình: một lời giải thích đơn giản hơn
  mô hình phải lệch khỏi nó ở ít nhất một đầu vào. Thêm điều kiện
  \tn{không suy biến}{non-degenerate}
  thì chỗ lệch ấy lớn hơn ngưỡng $\delta$, nên ô ở góc
  hình~\ref{fig:ch13-mien-kha-thi} rỗng. Định lý~\ref{thm:ch13-bo-ba-bat-kha},
  \tn{bộ ba bất khả}{trilemma} trong quy định, là đúng lập luận ấy viết bằng
  ngôn ngữ chính sách và không thêm toán học nào.
  \item Cái trần cao bao nhiêu thì tùy lớp hàm. Với hàm Boole ngẫu nhiên, mọi
  lời giải thích ngắn hơn gần trọn bảng chân trị đều sai một nửa số đầu vào. Với
  hàm $L$-Lipschitz, chặn trên tăng theo hàm mũ của số chiều và theo đa thức của
  nghịch đảo sai số, đó là lời nguyền số chiều đặt vào bài toán giải thích;
  nhưng bài chỉ chứng minh được một chặn dưới tăng tuyến tính theo số chiều, nên
  lời nguyền ấy là tính chất của phép dựng lưới chứ chưa phải của bài toán. Bài
  phát biểu rằng hai chặn khớp nhau tới hằng số và sách không nhận phát biểu ấy.
  Lối ra duy nhất trong bài là số chiều nội tại: nếu dữ liệu nằm
  quanh một cấu trúc chiều thấp thì số mũ trong chặn trên đổi theo chiều ấy.
  Bảng~\ref{tab:ch13-do-phuc-tap-lop-mo-hinh} là chỗ bài thay $K$ bằng phép đếm
  tham số cho từng lớp mô hình, và mỗi ô của nó là một chặn trên chứ không phải
  một giá trị. Sách cũng không dùng hệ quả 3.8 của bài, vì chứng minh của nó rút
  chặn dưới cho sai số ra từ một chặn trên cho độ phức tạp, cùng kiểu hụt ấy.
  \item Chỗ duy nhất bài chạm tới Phần~II là kết luận rằng lời giải thích cục bộ
  rẻ hơn lời giải thích toàn cục theo hàm mũ. Phép so ấy đặt một độ dài mô tả
  tương đối, đo khi chương trình được truy cập mô hình như một tiên tri, cạnh
  một độ dài mô tả tuyệt đối; mà phần được tiên tri cho không chính là số hạng
  trội của vế toàn cục. Bài phát biểu quy ước ấy trong chính định lý và không
  mang nó vào phép so.
  \item Cái trần nói được ba điều: có một cái sàn không phương pháp nào lách
  được, cái sàn ấy đến từ phép đếm chương trình ngắn, và sự đánh đổi có tốc độ
  đo được. Nó không nói gì về độ trung thực của định
  nghĩa~\ref{def:ch01-do-trung-thuc}, vì sai số của nó đo trên đầu ra và điểm
  tối ưu của nó, một chương trình tính đúng $f$, không có bảo đảm nào về cách mô
  hình đi tới đầu ra. Nó cũng không phải đáy của chuỗi kiểm chứng mà
  chương~\ref{ch:con-nguoi-vang-mat} để lại, vì nó không phải một dụng cụ đo.
  Cái nó đổi được là yêu cầu đặt lên một dụng cụ tương lai: dụng cụ ấy không thể
  có việc xác nhận trùng khớp, mà phải mô tả chỗ lệch.
\end{itemize}
\end{tomtat}

\cauhoi
\begin{cauhoids}
  \item Điều kiện không suy biến hỏng với mô hình hồi quy nhận giá trị thực. Hãy
  dựng một cặp $f$ và $g$ cụ thể trên $[0,1]$ sao cho $g(x) \ne f(x)$ tại mọi
  $x$ mà vẫn có $\mathcal{E}_{\infty}(f,g) \le \delta$ với một $\delta$ nhỏ cho
  trước, rồi cho biết trong bốn họ chỉ số của
  chương~\ref{ch:do-do-trung-thuc}, họ nào phát hiện được chỗ lệch ấy và họ nào
  không.
  \item Mục~\ref{sec:ch13-tran-noi-gi} lập luận rằng một chương trình tính đúng
  $f$ chưa chắc là một lời giải thích trung thực theo định
  nghĩa~\ref{def:ch01-do-trung-thuc}. Hãy nêu một điều kiện đặt lên chương trình
  ấy sao cho nếu thỏa thì nó trung thực, rồi cho biết điều kiện đó có phát biểu
  được trong khung của bài báo hay không, và vì sao.
  \item Định lý về độ phức tạp cục bộ đo độ dài mô tả khi chương trình được truy
  cập mô hình như một tiên tri. Hãy phát biểu lại chặn cho trường hợp không có
  tiên tri, tức tính cả chi phí mã hóa các giá trị hàm, rồi cho biết phép so với
  chặn toàn cục còn lại gì.
  \item Lối thoát theo số chiều nội tại giả định dữ liệu tụ quanh một cấu trúc
  chiều thấp. Hãy nêu hai cách mà lối thoát ấy hỏng trong thực hành, một cách
  quy về bản thân dữ liệu và một cách quy về việc sai số được đo dưới phân phối
  nào, rồi cho biết cách thứ hai liên hệ thế nào với phản đối về đầu vào ngoài
  phân phối ở mục~\ref{sec:ch08-huan-luyen-lai}.
  \item Hệ quả 3.8 của bài rút một chặn dưới cho sai số từ một chặn trên cho độ
  phức tạp. Hãy chỉ ra chính xác bước lập luận không đứng vững, rồi cho biết
  phải có thêm kết quả nào thì kết luận của hệ quả mới theo được.
  \item Định lý~\ref{thm:ch13-bo-ba-bat-kha} buộc một khung quy định bỏ một
  trong ba đòi hỏi. Hãy chọn một đòi hỏi, mô tả một quy định thật sẽ trông ra
  sao nếu bỏ nó, rồi cho biết chương nào trước đây trong sách cung cấp bằng
  chứng về cái giá phải trả cho lựa chọn ấy.
  \item Đọc mục 3.3 của \enquote{The Limits of AI Explainability: An Algorithmic
  Information Theory Approach} của Rao~\autocite{p26limits}, phần mà chương này
  không đọc. Hãy cho biết chặn theo thông tin tương hỗ ở đó mạnh hơn hay yếu hơn
  định lý~\ref{thm:ch13-khoang-cach-do-phuc-tap}, và nêu một giả thiết mà nó cần
  còn định lý~\ref{thm:ch13-khoang-cach-do-phuc-tap} thì không.
\end{cauhoids}
